\chapter{Beyond Standard Model (BSM) physics}
We now want to investigate the UV fields that could possibly leave their imprint on muon decay. The direction of our inquiry, as we are coming from the low-energy regime, makes the necessary line of reasoning quite intricate. Inferring the effective low-energy limit of a given UV extension of the SM is relatively straight-forward: We construct an amplitude with a heavy degree of freedom and integrate it out, as we did in Sec.~\ref{sec:intoutHDOF}. As we match the resulting expression to the corresponding effective contact interaction we find its Wilson coefficient. And as we carry out this calculation for all constructable amplitudes we derive the complete set of Wilson coefficients and thus the full low-energy EFT. However, finding all UV origins for a given effective operator by such a top-down approach requires extensive research -- which comes down to constructing all low-energy EFTs for all possible UV models (of a certain type). So for the directive of this work a more effective bottom-up strategy would clearly be more adequate.\par 
Luckily the bottom-up approach has been made effective by the work on local on-shell amplitudes of recent years. Specifically Li, Ni, Xiao \& Yu showed in \cite{li2022bottom} the correspondence of on-shell contact amplitudes and effective operators and De Angelis provided in \cite{de2022amplitude} an algorithm for finding all kinematically independent massive amplitudes from UV states. This allows Li, Ren \& Yu in \cite{li2024complete} to present a complete UV dictionary for SMEFT operators up to dimension seven. Precisely, their dictionary allows us to look up all BSM particles of spin $0,\tfrac{1}{2},1$ that lead to renormalizable UV interactions which would then match at tree-level to a SMEFT operator of our interest -- which for our endevaour means a SMEFT operator that could disturb muon decay.\par 
We order the steps of the thought process in the following way: First we look up all UV fields that produce the SMEFT operators we derived in Sec.~\ref{sec:matchingLEFTSMEFT}. In the following section we analyze the Lagrangian for a specific UV field -- $S_4$ -- so we are well prepared to derive our results in the next chapter.

\section{The UV dictionary}
\begin{table}[tbp]
    \centering
    %\renewcommand{\arraystretch}{1.5}
    \small
    \begin{tabular}{lccccccccc}
        \toprule
        Notation \cite{li2024complete}& $S_2$ & $S_4$ & $S_6$ & $F_1$ & $F_2$ & $F_3$&$F_4$\\
        Notation \cite{de2018effective} & $\mathcal{S}_1$ &$\varphi$ & $\Xi_1$ & $N$ & $E^c$ & $\Delta^c_1$& $\Delta^c_3$\\
        Irrep & $({\bf 1,1})_1$ & $({\bf 1,2})_{\frac{1}{2}}$ & $({\bf 1,3})_1$ & $({\bf 1,1})_0$ & $({\bf 1,1})_1$ & $({\bf 1,2})_{\frac{3}{2}}$&$({\bf 1,2})_{\frac{1}{2}}$\\
        \midrule
        Notation \cite{li2024complete}& $F_5$ & $F_6$ & $V_1$ & $V_2$ & $V_3$ & $V_4$&  \\
        Notation \cite{de2018effective} & $\Sigma$ & $\Sigma_1^c$ & $\mathcal{B}$ & $\mathcal{B}_1$ & $\mathcal{L}_3^\dagger$ & $\mathcal{W}$& \\
        Irrep & $({\bf 1,3})_0$& $({\bf 1,3})_1$& $({\bf 1,1})_0$ & $({\bf 1,1})_1$ & $({\bf 1,2})_{\frac{3}{2}}$ & $({\bf 1,3})_0$&\\
        \bottomrule
    \end{tabular}
    \caption{BSM scalars, fermions and vectors that produce the operators $\mathcal{O}_{ll}, \mathcal{O}_{le}, \mathcal{O}^{(1)}_{Hl}, \mathcal{O}^{(3)}_{Hl}, \mathcal{O}_{He},\mathcal{O}_{leHD}, \mathcal{O}_{\bar elllH}$ at tree-level. Fermions have a superscript ${}^c$ if they are conjugated with each other. Majorana fermions $F_1$ and $F_5$ have parity left, i.e. $F = F_L = P_L F$. Vectors have a superscript ${}^\dagger$ if they are conjugated with each other.}
    \label{tab:rBSMsfv}
\end{table}
We use the dictionary from \cite{li2024complete} to determine the possible UV origins of our SMEFT operators -- the complete list of new UV states that at tree-level produce SMEFT operators of dimension five, six and seven is cited in Tab.~\ref{tab:BSMsfv} in the appendix. We see that the dimension six operators are produced by adding a single new UV field to the Lagrangian. The four-lepton operators $\op_{ll}$ and $\op_{le}$ naturally need a bosonic mediator -- a vertex of three fermions would not produce a Lorentz scalar and have half-integer canonical mass dimension after all. The operators $\op^{(1,3)}_{Hl}$ on the other hand allows for both a vector and a fermion channel -- either connecting the currents $(H^\dagger DH)$ and $(ll)$ with each other or $H^\dagger l$ with $Hl$. Our two dimension seven operators $\mathcal{O}_{leHD}$ and $\mathcal{O}_{\bar elllH}$ are produced by either adding $F_1$ or $F_5$ or combinations of fields. The fermions $F_1\sim({\bf 1,1})_0$ and $F_5\sim({\bf 1,3})_0$ both have zero hypercharge and transform as real representations of $SU(2)_L$ -- so they are self-conjugate and their Majorana mass exactly provides the $\Delta L=2$ for producing a dimension seven SMEFT operator that violates lepton number by two. For the remaining completions, no single non-Majorana field suffices -- two heavy propagators are needed in order to attach all the external fields and violate lepton number. 
\begin{align}
    \mathcal{O}_{ll}&\to S_2,S_6, V_1, V_4\label{eq:bsmOll}\\
    \mathcal{O}_{le}&\to S_4, V_1, V_3\\
    \mathcal{O}^{(1)}_{Hl}&\to F_1,F_2,F_5,F_6,V_1\\
    \mathcal{O}^{(3)}_{Hl}&\to F_1,F_2,F_5,F_6,V_4\\
    \mathcal{O}_{He}&\to F_3,F_4,V_1\\
    \mathcal{O}_{leHD}&\to F_1,F_5,(F_3,S_6),(F_3,V_2),(V_2,V_3)\label{eq:bsmOleHD}\\
    \mathcal{O}_{\bar elllH}&\to F_1,F_5,(S_2,S_4),(S_2,F_4),(S_4,S_6),(S_6,F_4)\label{eq:bsmOelllH}
\end{align}
The representations of all the UV fields above are given in Tab.~\ref{tab:rBSMsfv}. The notation in the first line of the table is that of Li, Ren \& Yu in \cite{li2024complete}, the notation in the second line the more common names as cited in \cite{de2018effective}. We would now like to analyze the (possible) UV origins of the dimension six coefficients and obtain mass limits. As we see in Eq.~\ref{eq:bsmOll} the coefficient $C_{ll}$ could be produced by scalar mediator $S_2$ or $S_6$ or a vector mediator $V_1$ or~$V_4$.\par 
Let us take a look at the precise matching conditions from the dictionary:
\begin{align}
\begin{autobreak}
 {C_{ll}} =
-\frac{\mathcal{D}_{{llS_{2}}}^{{f_1f_2p_1}} \mathcal{D}_{{llS_{2}}}^{{f_4f_3p_1*}}}{2 {M_{S_{2}}^2}}
+\frac{\mathcal{D}_{{llS_{6}}}^{{f_1f_2p_1}} \mathcal{D}_{{llS_{6}}}^{{f_4f_3p_1*}}}{4 {M_{S_{6}}^2}}
-\frac{\mathcal{D}_{{ll^\dagger V_{1}}}^{{f_1f_3p_1}} \mathcal{D}_{{ll^\dagger V_{1}}}^{{f_2f_4p_1}}}{2 {M_{V_{1}}^2}}
+\frac{\mathcal{D}_{{ll^\dagger V_{4}}}^{{f_1f_3p_1}} \mathcal{D}_{{ll^\dagger V_{4}}}^{{f_2f_4p_1}}}{4 {M_{V_{4}}^2}}
-\frac{\mathcal{D}_{{ll^\dagger V_{4}}}^{{f_1f_4p_1}} \mathcal{D}_{{ll^\dagger V_{4}}}^{{f_2f_3p_1}}}{2 {M_{V_{4}}^2}}
+\frac{\mathcal{D}_{{llS_{2}}}^{{f_2f_1p_1}} \mathcal{D}_{{llS_{2}}}^{{f_4f_3p_1*}}}{2 {M_{S_{2}}^2}}
+\frac{\mathcal{D}_{{llS_{6}}}^{{f_2f_1p_1}} \mathcal{D}_{{llS_{6}}}^{{f_4f_3p_1*}}}{4 {M_{S_{6}}^2}} 
\end{autobreak}
\end{align}
The indices $(f_1,f_2,f_3,f_4)$ are flavor indices which in Eq.~\ref{Oll} are labeled $(s,r,t,p)$. We see that the matching conditions have the following general structure:
\begin{equation}
    C_i^{(6)}(M_X)=\kappa\frac{\lambda_a\lambda_b^*}{M_X^2}
\end{equation}
Here we use $i=ll,le,Hl(1),Hl(3),He$ (the full matching conditions are cited in Eq.~\ref{eq:ll-UV} to Eq.~\ref{eq:elllH-UV} in the appendix) and $a,b$ are flavor indices and $\kappa$ is a numerical or group theoretic factor from the matching procedure. If we denote by $R$ the correction from the (LEFT or SMEFT) RGE and remember that our constrain $\epsilon_\text{allowed}$ relates to the LEFT coefficient by a factor of VEV squared $\epsilon_\text{allowed}\equiv v^2L_\text{allowed}(\mu_\mu)$ we can derive the general formula for the mass bounds on our heavy mediators:
\begin{equation}
    M_X>v\cdot\sqrt{\frac{|R\kappa\lambda_a\lambda_b^*|}{\epsilon_\text{allowed}}}
\end{equation}
For the dimension seven operators we take a look at the matching condition for $C_{leHD}$ -- the operator $\op_{\bar e lllH}$ has quite an opulent matching and we only cite it in the appendix in Eq.~\ref{eq:elllH-UV}. 
\begin{align}
\begin{autobreak}
 {C_{leHD}} =
-\frac{{y}_{{e}}^{{p_1f_5}} \mathcal{D}_{{F_{1}Hl}}^{{p_2f_1}} \mathcal{D}_{{F_{1}Hl}}^{{p_2p_1}}}{2 {M_{F_{1}}^3}}
-\frac{{y}_{{e}}^{{p_1f_5}} \mathcal{D}_{{llS_{6}}}^{{f_1p_1p_2}} \mathcal{C}_{{HHS_{6}^\dagger }}^{{p_2}}}{{M_{S_{6}}^4}}
-\frac{{y}_{{e}}^{{p_1f_5}} \mathcal{D}_{{F_{5}Hl}}^{{p_2f_1}} \mathcal{D}_{{F_{5}Hl}}^{{p_2p_1}}}{4 {M_{F_{5}}^3}}
+\frac{{y}_{{e}}^{{p_1f_5}} \mathcal{C}_{{HHS_{6}^\dagger }}^{{p_2}} \mathcal{D}_{{llS_{6}}}^{{p_1f_1p_2}}}{{M_{S_{6}}^4}}
+\frac{2 i \mathcal{D}_{{eF_{1}V_{2}}}^{{f_5p_1p_2}} \mathcal{D}_{{F_{1}Hl}}^{{p_1f_1}} \mathcal{D}_{{HHV_{2}^\dagger D}}^{{p_2}}}{{M_{F_{1}}} {M_{V_{2}}^2}}
-\frac{\mathcal{D}_{{e^\dagger F_{3R}^\dagger H^\dagger }}^{{f_5p_1*}} \mathcal{D}_{{F_{1}F_{3R}^\dagger H}}^{{p_2p_1}} \mathcal{D}_{{F_{1}Hl}}^{{p_2f_1}}}{{M_{F_{1}}} {M_{F_{3}}^2}}
-\frac{\mathcal{D}_{{e^\dagger F_{3R}^\dagger H^\dagger }}^{{f_5p_1*}} \mathcal{D}_{{F_{3L}F_{5}H^\dagger }}^{{p_1p_2*}} \mathcal{D}_{{F_{5}Hl}}^{{p_2f_1}}}{{M_{F_{3}}} {M_{F_{5}}^2}}
+\frac{2 i \mathcal{D}_{{e^\dagger F_{3R}^\dagger H^\dagger }}^{{f_5p_1*}} \mathcal{D}_{{F_{3L}l^\dagger V_{2}^\dagger }}^{{p_1f_1p_2*}} \mathcal{D}_{{HHV_{2}^\dagger D}}^{{p_2}}}{{M_{F_{3}}} {M_{V_{2}}^2}}
-\frac{\mathcal{D}_{{e^\dagger F_{3R}^\dagger H^\dagger }}^{{f_5p_1*}} \mathcal{D}_{{F_{3R}^\dagger F_{5}H}}^{{p_1p_2}} \mathcal{D}_{{F_{5}Hl}}^{{p_2f_1}}}{2 {M_{F_{3}}^2} {M_{F_{5}}}}
-\frac{\mathcal{D}_{{e^\dagger F_{3R}^\dagger H^\dagger }}^{{f_5p_1*}} \mathcal{D}_{{F_{3R}^\dagger lS_{6}}}^{{p_1f_1p_2}} \mathcal{C}_{{HHS_{6}^\dagger }}^{{p_2}}}{{M_{F_{3}}^2} {M_{S_{6}}^2}}
+\frac{\mathcal{D}_{{e^\dagger F_{5}S_{6}^\dagger }}^{{f_5p_1p_2*}} \mathcal{D}_{{F_{5}Hl}}^{{p_1f_1}} \mathcal{C}_{{HHS_{6}^\dagger }}^{{p_2}}}{{M_{F_{5}}^2} {M_{S_{6}}^2}}
-\frac{2 i \mathcal{D}_{{e^\dagger l^\dagger V_{3}^\dagger }}^{{f_5f_1p_1*}} \mathcal{C}_{{HHS_{6}^\dagger }}^{{p_2}} \mathcal{D}_{{HS_{6}V_{3}^\dagger D}}^{{p_2p_1}}}{{M_{S_{6}}^2} {M_{V_{3}}^2}}
-\frac{4 i \mathcal{D}_{{e^\dagger l^\dagger V_{3}^\dagger }}^{{f_5f_1p_1*}} \mathcal{D}_{{HHV_{2}^\dagger D}}^{{p_2}} \mathcal{C}_{{HV_{2}V_{3}^\dagger }}^{{p_2p_1}}}{{M_{V_{2}}^2} {M_{V_{3}}^2}} 
\end{autobreak}\label{eq:leHD-UV}
\end{align}
We see that for a SM extension of a single fermion $F_1$ or $F_5$ the mass is bound by a term of the following form:
\begin{equation}
    M_X=v\cdot\left(\frac{|R\kappa\lambda_\text{eff}|}{\epsilon_\text{allowed}}\right)^\frac{1}{3}
\end{equation}
For a two-field completion the experiment only constraints:
\begin{equation}
    M_X^2M_Y>\frac{|R\kappa\lambda_1\lambda_2\lambda_3|v^3}{\epsilon_\text{allowed}}
\end{equation}
Only if we assume $M_X\approx M_Y$ the cube root formula applies.
\section{Correlated mediators}
Especially interesting are mediators that produce multiple SMEFT operators that disturb muon decay -- as for example the BSM vector field $V_1$, which matches onto $\op_{ll},\op_{le}$ and $\op_{Hl}^{(1)}$ all-together. Schematically $V_1$ couples to three currents:
\begin{equation}
    \mathcal{L}_\text{BSM}\supset V_{1\mu}(\lambda_HJ^\mu_H+\lambda_{l,pr}J_l^{\mu,pr}+\lambda_{e,pr}J_e^{\mu,pr})
\end{equation}
So integrating out $V_1$ generates simultaneously:
\begin{equation}
    C_{ll}\sim\frac{\lambda_l\lambda_l}{M_{V_1}^2}\qquad C_{le}\sim\frac{\lambda_l\lambda_e}{M_{V_1}^2}\qquad 
    C_{Hl}^{(1)}\sim\frac{\lambda_H\lambda_l}{M_{V_1}^2}\label{eq:V}
\end{equation}
These should be constrained together by the LEFT observables -- which we will do in the following chapter. 

